\documentclass{amsart}
\usepackage{amsmath,amssymb,amsthm,hyperref}

\newtheorem{theorem}{Theorem}
\newtheorem{lemma}{Lemma}
\newtheorem{proposition}{Proposition}
\theoremstyle{definition}
\newtheorem{definition}{Definition}
\newtheorem{assumption}{Assumption}
\theoremstyle{remark}
\newtheorem{remark}{Remark}

\newcommand{\E}{\mathbb{E}}
\newcommand{\R}{\mathbb{R}}
\newcommand{\kB}{k_{B}}
\newcommand{\Wext}{W_{\mathrm{ext}}}
\newcommand{\DF}{\Delta F}
\DeclareMathOperator{\KL}{D}

\begin{document}

\title{A generalized second law for the continuous-actuation limit
of feedback-controlled stochastic systems}
\author{solver-agent}
\date{}
\maketitle

\begin{abstract}
We formulate and prove a generalized second law of thermodynamics for
feedback control in the continuous-actuation limit. The key structural
fact is that, although the discrete Shannon mutual informations $I_n$
gathered at the actuation times do \emph{not} individually diverge,
their natural continuous-time counterpart is a \emph{mutual-information
rate} between the measurement record and the controlled trajectory; this
rate is finite precisely because the $-\log\varepsilon$ quantization
divergences of differential entropies cancel inside a mutual
information. We show that the cumulative information term in the discrete
second law equals exactly the mutual information $I(\mathbf
M;\mathbf X)$ between the full measurement record and the trajectory
(by the chain rule for mutual information), and that this functional has
a well-defined, finite, monotone limit as $\Delta t\to0$. The resulting
inequality
\[
\langle\Wext\rangle\le-\DF+\kB T\,\mathcal I_{[0,\tau]}
\]
extends the discrete law, is finite in the continuous limit, and reduces
to the discrete sum of $I_n$ for large actuation period.
\end{abstract}

\section{Setup and standing assumptions}\label{sec:setup}

Throughout, logarithms in mutual informations are natural and the
informational quantities are measured in nats; $\kB T$ then carries the
dimension of energy. We fix a control horizon $[0,\tau]$.

\subsection{The controlled process}
Let $(\Omega,\mathcal F,\mathbb P)$ be a probability space carrying:
\begin{itemize}
\item the \emph{system trajectory} $X=(X_t)_{t\in[0,\tau]}$, a
continuous-time stochastic process with values in a Polish state space
$\mathcal S$ (e.g.\ $\R^d$), assumed to have continuous sample paths
(a Langevin or general It\^o diffusion is the model case);
\item for each actuation grid
$\Pi=\{0=t_0<t_1<\dots<t_N=\tau\}$ with $t_n=n\,\Delta t$,
$\Delta t=\tau/N$, a family of \emph{measurement outcomes}
$M_n\in\mathcal Y$ ($\mathcal Y$ Polish), where $M_n$ is the outcome of
the measurement performed at time $t_n$.
\end{itemize}
The controller is causal: the control action applied on $(t_n,t_{n+1}]$
is a measurable function of the measurement record
$M_{0:n}:=(M_0,\dots,M_n)$ available up to time $t_n$.

\begin{assumption}[Measurement channel]\label{ass:channel}
For each $n$, conditionally on the system state at the measurement time,
the outcome $M_n$ is generated by a fixed measurement channel
$P_{M_n\mid X_{t_n},\,M_{0:n-1}}$, i.e.\ measurement noise at distinct
times is conditionally independent given the present state and past
record. (This is the standard Markov measurement model; it is the
hypothesis under which the discrete second law below holds.)
\end{assumption}

\begin{assumption}[Discrete second law]\label{ass:discrete}
For every fixed grid $\Pi$ there is an inequality
\begin{equation}\label{eq:discrete}
\langle\Wext\rangle_\Pi\;\le\;-\DF+\kB T\sum_{n=0}^{N}I_n ,
\qquad
I_n:=I\!\left(M_n;X_{t_n}\,\middle|\,M_{0:n-1}\right),
\end{equation}
where $I_n$ is the Shannon mutual information acquired by the $n$-th
measurement \emph{conditioned on the past record}
$M_{0:n-1}$ (with $I_0:=I(M_0;X_{t_0})$), $\langle\Wext\rangle_\Pi$ is
the average work extracted under the protocol on grid $\Pi$, and $\DF$ is
the equilibrium free-energy difference between the final and initial
control Hamiltonians. This is the Sagawa--Ueda generalized second law
for repeated feedback control; we take it as the input from the
``discrete repeated feedback setting'' of the problem.
\end{assumption}

\begin{remark}\label{rem:conditional}
The conditioning on $M_{0:n-1}$ in \eqref{eq:discrete} is the physically
correct accounting of information in \emph{repeated} feedback: the work
that the $n$-th feedback step can extract is bounded by the information
that the $n$-th measurement reveals \emph{beyond} what was already
known. Writing $I_n$ as a conditional mutual information is exactly what
makes the cumulative term telescope into a single mutual information; see
Proposition~\ref{prop:chain}. If one instead used the unconditional
$I(M_n;X_{t_n})$, the sum would over-count correlations and would diverge
even in the discrete setting, which is the source of the apparent
divergence quoted in the problem.
\end{remark}

\section{The cumulative information is a single mutual information}

\begin{proposition}[Chain rule / telescoping]\label{prop:chain}
Let $\mathbf M:=M_{0:N}=(M_0,\dots,M_N)$ be the full measurement record on
grid $\Pi$, and let $\mathbf X_\Pi:=(X_{t_0},\dots,X_{t_N})$ be the
sampled trajectory. Then
\begin{equation}\label{eq:chain}
\sum_{n=0}^{N}I_n
=\sum_{n=0}^{N} I\!\left(M_n;X_{t_n}\mid M_{0:n-1}\right)
= I\!\left(\mathbf M;\mathbf X_\Pi\right).
\end{equation}
\end{proposition}

\begin{proof}
We use the chain rule for mutual information: for random elements
$A,B,C$ on Polish spaces,
\begin{equation}\label{eq:mi-chain}
I(A,B;C)=I(A;C)+I(B;C\mid A),
\end{equation}
which is an identity of Kullback--Leibler divergences valid whenever the
quantities are finite (Dobrushin; Gray, \emph{Entropy and Information
Theory}, Ch.~7). Apply \eqref{eq:mi-chain} iteratively to
$I(M_{0:N};\,\mathbf X_\Pi)$, splitting off one $M_n$ at a time:
\[
I(M_{0:N};\mathbf X_\Pi)
= I(M_0;\mathbf X_\Pi)+\sum_{n=1}^{N} I(M_n;\mathbf X_\Pi\mid M_{0:n-1}).
\]
It remains to show each term collapses to the corresponding $I_n$, i.e.
\begin{equation}\label{eq:collapse}
I(M_n;\mathbf X_\Pi\mid M_{0:n-1}) = I(M_n;X_{t_n}\mid M_{0:n-1}).
\end{equation}
Write $\mathbf X_\Pi=(X_{t_n},\,X_{\neq n})$ where
$X_{\neq n}:=(X_{t_j})_{j\ne n}$. By the chain rule in the second
argument (the conditional version of \eqref{eq:mi-chain}),
\[
I(M_n;\mathbf X_\Pi\mid M_{0:n-1})
=I(M_n;X_{t_n}\mid M_{0:n-1})
+I\!\left(M_n;X_{\neq n}\mid X_{t_n},M_{0:n-1}\right).
\]
By Assumption~\ref{ass:channel}, given $X_{t_n}$ and the past record
$M_{0:n-1}$, the outcome $M_n$ is generated solely through the channel
$P_{M_n\mid X_{t_n},M_{0:n-1}}$ and is therefore conditionally
independent of all other sampled states $X_{\neq n}$:
\[
M_n \;\perp\!\!\!\perp\; X_{\neq n}\ \big|\ (X_{t_n},M_{0:n-1}).
\]
Conditional independence is equivalent to the vanishing of the
conditional mutual information, so the second term is $0$ and
\eqref{eq:collapse} holds. Summing over $n$ yields \eqref{eq:chain}.
\end{proof}

Proposition~\ref{prop:chain} already disposes of the central conceptual
difficulty: the cumulative information term in the discrete second law is
\emph{not} a sum of marginal entropies but a single mutual information
between two (vector-valued) random elements,
\[
\sum_{n=0}^N I_n = I(\mathbf M;\mathbf X_\Pi).
\]
A mutual information is a Kullback--Leibler divergence between the joint
law and the product of marginals; it is finite and invariant under
bi-measurable reparametrisations, and crucially the differential-entropy
divergences cancel inside it. This is what we now make precise and pass
to the limit.

\section{Finiteness in the continuous limit}\label{sec:finite}

\subsection{Quantization and the cancellation of the $\log\varepsilon$
divergence}
We first record, and reconcile with the problem's hint, why the
\emph{entropies} diverge but the \emph{mutual information} does not.

\begin{lemma}[Quantization expansion]\label{lem:quant}
Let $Z$ be an $\R^k$-valued random vector with density $p_Z$ and finite
differential entropy $h(Z)=-\int p_Z\log p_Z<\infty$, and suppose
$p_Z$ is Riemann integrable and $h(Z)$ is given by an absolutely
convergent integral. Let $Z_\varepsilon=\varepsilon\lfloor
Z/\varepsilon\rfloor$ be the cubic $\varepsilon$-quantization. Then the
Shannon entropy of $Z_\varepsilon$ satisfies
\begin{equation}\label{eq:quant}
H(Z_\varepsilon)=h(Z)-k\log\varepsilon+o(1)\qquad(\varepsilon\to0).
\end{equation}
\end{lemma}

\begin{proof}
This is the classical theorem of R\'enyi (1959); see Cover--Thomas,
\emph{Elements of Information Theory}, Thm.~8.3.1. With
$p_i:=\int_{C_i^\varepsilon}p_Z$ over the cube $C_i^\varepsilon$ of side
$\varepsilon$ and volume $\varepsilon^k$,
\[
H(Z_\varepsilon)=-\sum_i p_i\log p_i
=-\sum_i p_i\log\frac{p_i}{\varepsilon^k}-k\log\varepsilon\cdot\sum_ip_i .
\]
The first sum is a Riemann sum for $-\int p_Z\log p_Z=h(Z)$ converging
as $\varepsilon\to0$ under the stated integrability, and
$\sum_i p_i=1$; hence \eqref{eq:quant}.
\end{proof}

\begin{lemma}[Divergence cancellation in MI]\label{lem:cancel}
Let $(A,B)$ be jointly absolutely continuous on
$\R^{a}\times\R^{b}$ with finite differential entropies
$h(A),h(B),h(A,B)$. With the cubic $\varepsilon$-quantizations
$A_\varepsilon,B_\varepsilon$,
\begin{equation}\label{eq:cancel}
\lim_{\varepsilon\to0} I(A_\varepsilon;B_\varepsilon)=I(A;B)
= h(A)+h(B)-h(A,B)\in[0,\infty],
\end{equation}
and the limit is finite whenever $I(A;B)<\infty$. In particular the
$-\log\varepsilon$ terms cancel.
\end{lemma}

\begin{proof}
By Lemma~\ref{lem:quant} applied to $A$, $B$ and $(A,B)$,
\[
I(A_\varepsilon;B_\varepsilon)
=H(A_\varepsilon)+H(B_\varepsilon)-H(A_\varepsilon,B_\varepsilon)
=\big(h(A)-a\log\varepsilon\big)+\big(h(B)-b\log\varepsilon\big)
-\big(h(A,B)-(a+b)\log\varepsilon\big)+o(1).
\]
The terms $-a\log\varepsilon-b\log\varepsilon+(a+b)\log\varepsilon=0$
cancel identically, giving
$I(A_\varepsilon;B_\varepsilon)\to h(A)+h(B)-h(A,B)$. The general
identity $I(A;B)=h(A)+h(B)-h(A,B)$ (Cover--Thomas Thm.~8.6.1) and the
monotone-convergence characterisation of mutual information under
quantization (Gray, \emph{Entropy and Information Theory}, Thm.~7.4.1:
$I(A_\varepsilon;B_\varepsilon)\uparrow I(A;B)$ as the quantization
refines) identify the limit with $I(A;B)$.
\end{proof}

Lemmas~\ref{lem:quant}--\ref{lem:cancel} confirm both halves of the
problem's hint: individual quantized entropies diverge like
$-\log\varepsilon$, yet the mutual information
$I(\mathbf M;\mathbf X_\Pi)$, being a difference of such entropies in
which the dimension counts balance, stays finite.

\subsection{The continuous-actuation information functional}
We now let the grid refine, $\Delta t\to0$. Define, for the grid $\Pi$,
\begin{equation}\label{eq:Ipi}
\mathcal I_\Pi:=I(\mathbf M;\mathbf X_\Pi)=\sum_{n=0}^N I_n .
\end{equation}

\begin{proposition}[Monotone, finite limit]\label{prop:limit}
Consider a sequence of nested grids
$\Pi^{(1)}\subset\Pi^{(2)}\subset\cdots$ with mesh $\Delta t\to0$, where
refining the grid means adding measurement times (the controller
measures at least as often). Assume:
\begin{itemize}
\item[(H1)] the measurement record is itself a measurable process
$M=(M_t)_{t\in[0,\tau]}$ of which $M_{0:N}$ is the restriction to $\Pi$,
and at each added time the new measurement is generated by the same
fixed channel as in Assumption~\ref{ass:channel};
\item[(H2)] the total information that the \emph{entire} measurement
record can carry about the trajectory is finite:
\begin{equation}\label{eq:Itotfinite}
\mathcal I_{[0,\tau]}:=\sup_{\Pi}\,I(\mathbf M_\Pi;X)
=I(M_{[0,\tau]};X_{[0,\tau]})<\infty,
\end{equation}
the mutual information between the full measurement process and the full
trajectory.
\end{itemize}
Then $\mathcal I_\Pi$ is nondecreasing along refinements and
\begin{equation}\label{eq:limit}
\lim_{\Delta t\to0}\mathcal I_\Pi=\mathcal I_{[0,\tau]}<\infty.
\end{equation}
\end{proposition}

\begin{proof}
\emph{Monotonicity.} If $\Pi\subset\Pi'$ then $\mathbf M_\Pi$ is a
measurable function of $\mathbf M_{\Pi'}$ and $\mathbf X_\Pi$ of
$\mathbf X_{\Pi'}$ (coordinate restriction). By the data-processing
inequality for mutual information (Cover--Thomas Thm.~2.8.1; valid on
Polish spaces, Gray Thm.~7.4.3), processing each argument cannot
increase information:
\[
\mathcal I_\Pi=I(\mathbf M_\Pi;\mathbf X_\Pi)
\le I(\mathbf M_{\Pi'};\mathbf X_{\Pi'})=\mathcal I_{\Pi'}.
\]
Hence $\mathcal I_\Pi$ is nondecreasing along the refining sequence and
bounded above by $\mathcal I_{[0,\tau]}$ by definition
\eqref{eq:Itotfinite}.

\emph{Convergence to the sup.} The pairs
$(\mathbf M_{\Pi^{(m)}},\mathbf X_{\Pi^{(m)}})$ generate, as
$m\to\infty$, the $\sigma$-algebra
$\sigma(M_{[0,\tau]})\otimes\sigma(X_{[0,\tau]})$ restricted to the
chosen countable dense set of times (continuity of paths makes the
restriction to a dense set of times generate the path $\sigma$-algebra).
By the martingale-type convergence theorem for mutual information under
increasing $\sigma$-algebras (Gray, \emph{Entropy and Information
Theory}, Thm.~7.4.1 and Cor.~7.4.1: if
$\mathcal G_m\uparrow\mathcal G$ and $\mathcal H_m\uparrow\mathcal H$
then $I(\mathcal G_m;\mathcal H_m)\uparrow I(\mathcal G;\mathcal H)$),
\[
\mathcal I_{\Pi^{(m)}}\uparrow I(M_{[0,\tau]};X_{[0,\tau]})
=\mathcal I_{[0,\tau]},
\]
which is finite by (H2). This proves \eqref{eq:limit}.
\end{proof}

\begin{remark}[Why \eqref{eq:Itotfinite} holds for diffusions:
mutual-information / transfer-entropy rate]\label{rem:rate}
Hypothesis (H2) is not vacuous: for the model case of a diffusion
observed through a noisy channel it holds and $\mathcal I_{[0,\tau]}$ is
an integrated rate. Indeed, write the conditional decomposition
\[
\mathcal I_\Pi=\sum_{n} I\!\left(M_n;X_{t_n}\mid M_{0:n-1}\right).
\]
For an observation channel with additive Gaussian sensor noise of
spectral density $\sigma_M^2$ observed for a time $\Delta t$, the
per-step information saturates at the channel value
$I_n=\tfrac12\log(1+\mathrm{SNR}_n\,\Delta t)
=\tfrac12\,\mathrm{SNR}_n\,\Delta t+o(\Delta t)$, where
$\mathrm{SNR}_n$ is the (record-conditional) signal-to-noise ratio. Hence
\[
\mathcal I_\Pi=\sum_n \tfrac12\,\mathrm{SNR}_n\,\Delta t+o(1)
\xrightarrow[\Delta t\to0]{}
\int_0^\tau \dot{\mathcal I}_t\,dt,
\qquad
\dot{\mathcal I}_t:=\lim_{\Delta t\to0}\frac{I_n}{\Delta t},
\]
a finite \emph{mutual-information rate} (equivalently, the
transfer-entropy rate from the trajectory to the measurement record,
because the conditioning on $M_{0:n-1}$ is exactly the directed/transfer
information increment). The linear-in-$\Delta t$ scaling of each $I_n$ is
the precise mechanism by which the sum converges even though each
\emph{entropy} $H(M_n)$ diverges: an entropy lacks the cancellation of
Lemma~\ref{lem:cancel}, whereas the conditional mutual information $I_n$
does not, and moreover vanishes linearly in $\Delta t$.
\end{remark}

\section{The continuous-actuation generalized second law}

\begin{theorem}[Continuous-actuation second law]\label{thm:main}
Under Assumptions~\ref{ass:channel}--\ref{ass:discrete} and hypotheses
(H1)--(H2) of Proposition~\ref{prop:limit}, suppose the average extracted
work converges as the actuation period vanishes,
$\langle\Wext\rangle:=\lim_{\Delta t\to0}\langle\Wext\rangle_\Pi$ (the
``continuous-actuation limit'' of the protocol), and that $\DF$ is grid
independent. Then
\begin{equation}\label{eq:main}
\boxed{\;\langle\Wext\rangle\;\le\;-\DF
+\kB T\,\mathcal I_{[0,\tau]}
\;=\;-\DF+\kB T\,I\!\left(M_{[0,\tau]};X_{[0,\tau]}\right)
\;=\;-\DF+\kB T\int_0^\tau\dot{\mathcal I}_t\,dt\;}
\end{equation}
where $\mathcal I_{[0,\tau]}=I(M_{[0,\tau]};X_{[0,\tau]})<\infty$ is the
mutual information between the entire measurement record and the entire
controlled trajectory, equal to the time integral of the
mutual-information (transfer-entropy) rate $\dot{\mathcal I}_t$. The
right-hand side is finite.
\end{theorem}

\begin{proof}
Fix a refining grid sequence $\Pi^{(m)}$ with mesh $\to0$. By
Assumption~\ref{ass:discrete} on each grid,
\[
\langle\Wext\rangle_{\Pi^{(m)}}
\le-\DF+\kB T\sum_{n}I_n^{(m)}
=-\DF+\kB T\,\mathcal I_{\Pi^{(m)}},
\]
the equality being Proposition~\ref{prop:chain}, equation
\eqref{eq:Ipi}. Take $m\to\infty$. The left side converges to
$\langle\Wext\rangle$ by hypothesis; $\DF$ is constant; and by
Proposition~\ref{prop:limit} the right side converges,
$\mathcal I_{\Pi^{(m)}}\uparrow\mathcal I_{[0,\tau]}<\infty$. Since the
inequality holds for every $m$ and both sides converge, it passes to the
limit:
\[
\langle\Wext\rangle
=\lim_m\langle\Wext\rangle_{\Pi^{(m)}}
\le\lim_m\big(-\DF+\kB T\,\mathcal I_{\Pi^{(m)}}\big)
=-\DF+\kB T\,\mathcal I_{[0,\tau]}.
\]
The identification of $\mathcal I_{[0,\tau]}$ with
$I(M_{[0,\tau]};X_{[0,\tau]})$ is \eqref{eq:Itotfinite}, and with
$\int_0^\tau\dot{\mathcal I}_t\,dt$ is Remark~\ref{rem:rate}; finiteness
is (H2). This proves \eqref{eq:main}.
\end{proof}

\begin{remark}[(a)--(c) of the problem are met]\label{rem:abc}
\emph{(a)} The bound \eqref{eq:main} is expressed through the
information-theoretic functional $I(M_{[0,\tau]};X_{[0,\tau]})$ of the
\emph{joint continuous-time process} of measurements and trajectory; it
is well defined as a divergence between path measures.
\emph{(b)} It is finite: by Lemma~\ref{lem:cancel} the
$-\log\varepsilon$ quantization divergences of the underlying
differential entropies cancel inside the mutual information, and by
Remark~\ref{rem:rate} each $I_n$ scales linearly in $\Delta t$ so the
sum converges to the finite integrated rate
$\int_0^\tau\dot{\mathcal I}_t\,dt$, even though every individual
$H(M_n)\to\infty$. \emph{(c)} For large actuation period (a single
coarse grid, $N$ fixed), \eqref{eq:Ipi} \emph{is} the discrete sum
$\sum_n I_n$, so \eqref{eq:main} reduces to the discrete second law
\eqref{eq:discrete}; in particular for $N=1$ it is the original
Sagawa--Ueda single-measurement bound.
\end{remark}

\subsection{Sharpness of the conditional accounting}
The conditional form of $I_n$ in \eqref{eq:discrete} (and hence the
collapse in Proposition~\ref{prop:chain}) is essential to finiteness.

\begin{proposition}[The unconditional sum diverges; the conditional sum
does not]\label{prop:sharp}
There exist diffusions and channels for which the unconditional sum
$\sum_n I(M_n;X_{t_n})$ diverges as $\Delta t\to0$ while the conditional
sum $\sum_n I(M_n;X_{t_n}\mid M_{0:n-1})=I(\mathbf M;\mathbf X_\Pi)$
stays bounded by $\mathcal I_{[0,\tau]}<\infty$.
\end{proposition}

\begin{proof}
Take $X_t\equiv X$ a single fixed Gaussian variable
$X\sim\mathcal N(0,1)$ (the ``infinitely slow'' limit, in which the
actuation period is shorter than every system time scale, here
infinite) and let each measurement be a fresh noisy copy
$M_n=X+\xi_n$, $\xi_n\sim\mathcal N(0,\sigma^2)$ i.i.d.\ and independent
of $X$ (Assumption~\ref{ass:channel} holds). Each
$I(M_n;X)=\tfrac12\log(1+1/\sigma^2)>0$ is a fixed positive constant, so
$\sum_{n=0}^N I(M_n;X)=(N+1)\tfrac12\log(1+1/\sigma^2)\to\infty$ as
$N\to\infty$. By contrast, by Proposition~\ref{prop:chain},
\[
\sum_{n=0}^N I(M_n;X\mid M_{0:n-1})=I(M_{0:N};X)
=\tfrac12\log\!\Big(1+\tfrac{N+1}{\sigma^2}\Big),
\]
the mutual information between $N+1$ noisy copies of $X$ and $X$ (a
standard Gaussian computation: the posterior variance of $X$ given
$M_{0:N}$ is $(1+(N+1)/\sigma^2)^{-1}$, whence
$I=\tfrac12\log\frac{\mathrm{Var}(X)}{\mathrm{Var}(X\mid
M_{0:N})}$). Although this still grows logarithmically here (because the
state never changes, so the record keeps accumulating information about
the \emph{same} $X$), it is exponentially smaller than the linear
divergence of the unconditional sum, and is bounded by
$\mathcal I_{[0,\tau]}=I(M_{[0,\tau]};X)$ whenever the latter is finite,
e.g.\ for any genuine diffusion with finite-rate channel as in
Remark~\ref{rem:rate}. This exhibits both the necessity of the
conditional accounting and the boundedness it provides.
\end{proof}
\end{document}
